%% mathstc-suites-lin-vs-geo — croissance linéaire (arithmétique, +60 par an) contre
%% croissance exponentielle (géométrique, x 1,12 par an). Le modèle exponentiel part
%% moins vite, puis rattrape et dépasse le modèle linéaire entre n = 4 et n = 5.
%% Échelles : 1 année -> 0,78 cm ; 100 unités de production -> 0,45 cm.
\documentclass[tikz,border=4pt]{standalone}
\usepackage{liaisons}
\begin{document}
\begin{tikzpicture}[x=1cm,y=1cm,
  axe/.style={-{Stealth[length=5pt]}, line width=0.6pt},
  serieA/.style={solideA, line width=0.7pt},
  serieB/.style={solideB, line width=0.7pt},
  grille/.style={line width=0.3pt, black!15},
  grad/.style={font=\scriptsize, inner sep=2.5pt},
  note/.style={font=\scriptsize, text=black!60, inner sep=1.5pt}]

\newcommand\xn[1]{{(#1)*0.78}}        % abscisse d'une année
\newcommand\yp[1]{{(#1)*0.0045}}      % ordonnée d'une production

%% ================= bande de croisement ====================================
\fill[black!60, opacity=0.15] (\xn{4},0) rectangle (\xn{5},\yp{1040});

%% ================= quadrillage horizontal =================================
\foreach \v in {200,400,...,1000} { \draw[grille] (0,\yp{\v}) -- (\xn{8.4},\yp{\v}); }

%% ================= axes et graduations ====================================
\draw[axe] (-0.12,0) -- (\xn{9},0);
\node[grad, anchor=south east] at (\xn{9},0.05) {année $n$};
\draw[axe] (0,-0.12) -- (0,\yp{1140});
\node[grad, anchor=north west, xshift=2pt] at (0.02,\yp{1140}) {production};
\node[grad, anchor=north east] at (-0.02,-0.02) {$O$};
\foreach \n in {0,1,...,8} {
  \draw[line width=0.5pt] (\xn{\n},0) -- (\xn{\n},-0.08);
  \node[grad, anchor=north] at (\xn{\n},-0.08) {\n}; }
\foreach \v in {200,400,600,800,1000} {
  \draw[line width=0.5pt] (0,\yp{\v}) -- (-0.08,\yp{\v});
  \node[grad, anchor=east] at (-0.08,\yp{\v}) {\v}; }

%% ================= série linéaire : 400 + 60n (carrés bleus) ==============
\draw[serieB] (\xn{0},\yp{400}) \foreach \n/\v in {1/460, 2/520, 3/580, 4/640,
                       5/700, 6/760, 7/820, 8/880} { -- (\xn{\n},\yp{\v}) };
\foreach \n/\v in {0/400, 1/460, 2/520, 3/580, 4/640, 5/700, 6/760, 7/820, 8/880} {
  \fill[solideB] ($(\xn{\n},\yp{\v})+(-1.5pt,-1.5pt)$) rectangle ++(3pt,3pt); }
\node[font=\scriptsize, text=solideB, anchor=west, inner sep=4pt]
     at (\xn{8},\yp{880}) {$400+60n$};

%% ================= série géométrique : 400 x 1,12^n (disques rouges) ======
\draw[serieA] (\xn{0},\yp{400}) \foreach \n/\v in {1/448, 2/501.8, 3/562.0, 4/629.4,
                       5/704.9, 6/789.5, 7/884.3, 8/990.4} { -- (\xn{\n},\yp{\v}) };
\foreach \n/\v in {0/400, 1/448, 2/501.8, 3/562.0, 4/629.4, 5/704.9, 6/789.5, 7/884.3, 8/990.4} {
  \fill[solideA] (\xn{\n},\yp{\v}) circle (1.8pt); }
\node[font=\scriptsize, text=solideA, anchor=west, inner sep=4pt]
     at (\xn{8},\yp{990.4}) {$400\times1{,}12^{\,n}$};

%% ================= les deux valeurs de part et d'autre du croisement ======
\node[grad, text=solideB, fill=white, inner sep=1pt] at ($(\xn{3.9},\yp{640})+(0,0.36cm)$) {$640$};
\node[grad, text=solideA, fill=white, inner sep=1pt] at ($(\xn{3.9},\yp{629.4})+(0,-0.36cm)$) {$629{,}4$};
\node[grad, text=solideA, fill=white, inner sep=1pt] at ($(\xn{5.15},\yp{704.9})+(0,0.36cm)$) {$704{,}9$};
\node[grad, text=solideB, fill=white, inner sep=1pt] at ($(\xn{5.15},\yp{700})+(0,-0.36cm)$) {$700$};

%% ================= légende de la bande ====================================
\node[note, anchor=north, align=center] at (\xn{4.5},-0.34)
     {le modèle exponentiel passe devant};
\end{tikzpicture}
\end{document}
